\chapter{LSZ Reduction}
\label{ch:lsz-reduction}

\section{The Reduction Problem}

Haag--Ruelle theory defines incoming and outgoing scattering states by
large-time Hilbert-space limits. LSZ reduction is the theorem that computes
matrix elements of that already-defined scattering operator from the singular
part of time-ordered Green functions at the one-particle mass shell.

This chapter keeps the same massive scalar particle as in the Haag--Ruelle
chapter. The assumptions are:
\begin{enumerate}
  \item the particle has mass \(m>0\) and an isolated one-particle subspace
  \(\Hilb_1\);
  \item the incoming and outgoing scattering subspaces for this particle
  coincide, so the scattering operator \(S=\Omega_{\mathrm{out}}^*
  \Omega_{\mathrm{in}}\) is unitary on the asymptotic Fock space;
  \item the scalar local field \(\widehat\phi\) has nonzero one-particle
  matrix element;
  \item the time-ordered products used below have been defined as
  distributions.
\end{enumerate}

Let generalized one-particle momentum states be normalized by
\[
  \langle \vec q\,|\vec p\rangle
  =
  (2\pi)^d\,2\omega_{\vec p}\,
  \delta^d(\vec q-\vec p),
  \qquad
  \omega_{\vec p}=\sqrt{\vec p^{\,2}+m^2},
\]
where \(d=D-1\). The Lorentz-invariant mass-shell measure is
\[
  \dd\mu_m(\vec p)
  =
  \frac{\dd^d\vec p}{(2\pi)^d\,2\omega_{\vec p}} .
\]
Kernel notation such as
\[
  {}_{\mathrm{out}}\langle \vec q_1,\ldots,\vec q_m
  |\vec p_1,\ldots,\vec p_n\rangle_{\mathrm{in}}
\]
means the distributional kernel obtained after smearing all momenta with
smooth compactly supported wave packets.

\section{The One-Particle Pole}

Assume that \(\widehat\phi\) is a scalar field with
\[
  \langle\Omega|\widehat\phi(0)|\vec p\rangle
  =
  Z^{1/2},
  \qquad Z>0,
\]
for the particle species under consideration. With the normalization above,
the K\"all\'en--Lehmann measure of \(\widehat\phi\) contains the atom
\[
  \dd\rho(s)=Z\,\delta(s-m^2)\,\dd s+\dd\rho_{\mathrm{cont}}(s).
\]
Consequently the Fourier transform of the time-ordered two-point function has
the pole
\[
  \widetilde G_2(k)
  =
  \frac{-\ii Z}{k^2+m^2-\ii0}
  +
  \widetilde G_{2,\mathrm{less}}(k),
\]
where \(\widetilde G_{2,\mathrm{less}}\) has no pole at the isolated
one-particle mass shell. The number \(Z\) depends on the chosen local field.
The particle state and the \(S\)-operator are invariant data; changing the
interpolating field changes the residue factor appearing in the reduction
formula.

The operation
\[
  Z^{-1/2}\,\ii(k^2+m^2)
\]
removes an external one-particle propagator with the conventions of the
Lorentzian Green-function chapter:
\[
  Z^{-1/2}\,\ii(k^2+m^2)\,
  \frac{-\ii Z}{k^2+m^2-\ii0}
  \longrightarrow
  Z^{1/2}
\]
as a mass-shell residue. Applying one such operation to every external field
is the analytic core of LSZ reduction.

\section{Connected Truncation and External Amputation}

The pole extracted by LSZ is an external one-particle pole.  It is useful to
separate this operation from two other operations that are often compressed
in notation.  First, the connected time-ordered functions are obtained from
the generating functional by taking cumulants.  If
\[
  Z[J]=
  \left\langle\Omega\left|
  \operatorname T\exp\left(\ii\int J(x)\widehat\phi(x)\,\dd^Dx\right)
  \right|\Omega\right\rangle,
\]
then
\[
  G_N^{\rm conn}(x_1,\ldots,x_N)
  =
  \left.
  \ii^{-N}
  \frac{\delta^N \log Z[J]}
       {\delta J(x_1)\cdots\delta J(x_N)}
  \right|_{J=0}.
\]
With the sign convention for \(J\) inherited from the exponential above, this
is the cumulant part of the time-ordered distribution: it contains precisely
the vacuum-connected components.  External one-particle propagators remain
until the residue operation is applied.

Second, amputation in LSZ is residue extraction at the isolated mass shell:
\[
  \operatorname*{Res}_{k^2=-m^2}
  \widetilde G_2(k)
  =
  -\ii Z .
\]
Multiplication by \(Z^{-1/2}\ii(k^2+m^2)\) is meaningful after smearing in a
neighborhood of the mass shell and taking the boundary value.  It removes
only the external pole associated with the chosen stable particle.  Internal
propagators, threshold singularities, bound-state poles in other channels,
and contact terms remain part of the amputated connected kernel.

Thus the reduction map has three logically distinct inputs:
\[
  \begin{gathered}
  \text{time-ordered distributions}
  \longrightarrow
  \text{connected distributions}
  \\
  \longrightarrow
  \text{external one-particle residues}
  \longrightarrow
  \text{scattering matrix elements}.
  \end{gathered}
\]
The Haag--Ruelle construction supplies the last target independently of
perturbation theory; the Green functions supply a way to compute its matrix
elements when the one-particle pole hypotheses hold.

\section{Contact Terms and Interpolating Fields}

The LSZ operation depends only on the simultaneous external one-particle pole.
This fact gives two useful stability properties.

First, different definitions of time-ordered products may differ by contact
terms supported on partial diagonals, for example on \(x_i=x_j\).  In momentum
space such terms are polynomial, or more generally local, in the momenta
associated with the coincident insertions.  They have no one-particle pole in
each external variable.  After multiplication by
\[
  \prod_a (k_a^2+m^2)
\]
and taking the on-shell boundary value with separated wave packets, these
local terms do not contribute to the connected scattering kernel.  They still
matter for Ward identities, composite operators, and renormalized local
products; they are simply not external one-particle residues.

Second, the local field used in the reduction is an interpolating field.  Let
\(\widehat\phi'\) be another scalar local field with
\[
  \langle\Omega|\widehat\phi'(0)|\vec p\rangle
  =
  a\,Z^{1/2},
  \qquad a\ne0,
\]
on the same isolated one-particle subspace.  Its two-point function has
one-particle residue \(Z'=|a|^2Z\).  The LSZ factors \(Z'^{-1/2}\) remove the
changed residue, and the resulting \(S\)-matrix element is the same.  Terms in
\(\widehat\phi'\) whose vacuum-to-one-particle matrix elements vanish, such
as equation-of-motion terms of the form \((\Box-m^2)\widehat\chi\) in these
sign conventions, do not create the external particle pole and therefore drop
out of the reduction formula.  The scattering operator belongs to the
asymptotic Hilbert space; the interpolating local fields are coordinates used
to access its matrix elements.

\section{Wave-Packet Form of LSZ}

Let
\[
  q_i=(\omega_{\vec q_i},\vec q_i),
  \qquad
  p_j=(\omega_{\vec p_j},\vec p_j)
\]
be outgoing and incoming positive-energy momenta. For test wave packets
\(f_i,g_j\in C_c^\infty(\mathbb R^d)\), define the smeared scattering matrix
element
\[
\begin{aligned}
&S[f_1,\ldots,f_m;g_1,\ldots,g_n]
\\
&\quad =
  \int
  \prod_{i=1}^m \dd\mu_m(\vec q_i)\, f_i(\vec q_i)^*
  \prod_{j=1}^n \dd\mu_m(\vec p_j)\, g_j(\vec p_j)
\\
&\qquad\qquad\times
  {}_{\mathrm{out}}\langle
  \vec q_1,\ldots,\vec q_m
  |\vec p_1,\ldots,\vec p_n
  \rangle_{\mathrm{in}} .
\end{aligned}
\]
The Haag--Ruelle construction expresses this number as a large-time limit of
vacuum matrix elements of local fields. Locality moves the fields associated
with future packets to the future side of the time-ordered product and the
fields associated with past packets to the past side. After Fourier transform,
the large-time oscillatory integrals select the one-particle poles
\[
  q_i^2=-m^2,
  \qquad
  p_j^2=-m^2.
\]

In distribution form, LSZ says that the connected part of the smeared matrix
element is obtained by replacing the scattering kernel with
\[
\begin{aligned}
&\left[
  \prod_{i=1}^m Z^{-1/2}\,\ii(q_i^2+m^2)
  \prod_{j=1}^n Z^{-1/2}\,\ii(p_j^2+m^2)
  \right]
\\
&\qquad\qquad\times
  \widetilde G^{\mathrm{conn}}_{m+n}
  (q_1,\ldots,q_m,-p_1,\ldots,-p_n)
\end{aligned}
\]
and then taking the boundary value in which every \(q_i\) and \(p_j\) lies on
\(\Sigma_m^+\). Incoming physical momenta appear as \(-p_j\) in the Green
function because all momenta in
\(\widetilde G_{m+n}\) are taken to enter the time-ordered correlator.

\section{Momentum-Space Kernel}

The time-ordered \(N\)-point function is
\[
  G_N(x_1,\ldots,x_N)
  =
  \langle\Omega|
  \operatorname T\{
    \widehat\phi(x_1)\cdots\widehat\phi(x_N)
  \}
  |\Omega\rangle .
\]
Its Fourier transform is defined by
\[
  G_N(x_1,\ldots,x_N)
  =
  \int
  \prod_{a=1}^N\frac{\dd^D k_a}{(2\pi)^D}\,
  \ee^{\ii\sum_a k_a\cdot x_a}\,
  \widetilde G_N(k_1,\ldots,k_N).
\]
Translation invariance gives
\[
  \widetilde G_N(k_1,\ldots,k_N)
  =
  (2\pi)^D\delta^D(k_1+\cdots+k_N)\,
  \widetilde{\mathcal G}_N(k_1,\ldots,k_{N-1}),
\]
where \(\widetilde{\mathcal G}_N\) denotes the reduced distribution after
removing the overall momentum-conservation delta function.

For \(m\) outgoing and \(n\) incoming scalar particles, set
\[
  k_i=q_i\quad (1\le i\le m),
  \qquad
  k_{m+j}=-p_j\quad (1\le j\le n).
\]
The connected LSZ kernel is
\[
\boxed{
\begin{aligned}
&{}_{\mathrm{out}}\langle
  \vec q_1,\ldots,\vec q_m
  |\vec p_1,\ldots,\vec p_n
  \rangle_{\mathrm{in,conn}}
\\
&\quad =
  \left[
  \prod_{a=1}^{m+n} Z^{-1/2}\,\ii(k_a^2+m^2)
  \right]
  \widetilde G^{\mathrm{conn}}_{m+n}(k_1,\ldots,k_{m+n})
  \bigg|_{\mathrm{on\ shell}} .
\end{aligned}
}
\]
The vertical bar means the distributional on-shell boundary value obtained
after smearing with wave packets. The formula is a statement about
distributions; the displayed kernel is the compact notation for that
statement.

Disconnected terms are treated by the same rule, but one-particle two-point
factors reproduce the identity contribution to the scattering operator. For
two identical scalar particles,
\[
\begin{aligned}
&{}_{\mathrm{out}}\langle
  \vec q_1,\vec q_2
  |\vec p_1,\vec p_2
  \rangle_{\mathrm{in}}
\\
&\quad =
  (2\pi)^d2\omega_{\vec p_1}\delta^d(\vec q_1-\vec p_1)\,
  (2\pi)^d2\omega_{\vec p_2}\delta^d(\vec q_2-\vec p_2)
\\
&\qquad+
  (2\pi)^d2\omega_{\vec p_1}\delta^d(\vec q_2-\vec p_1)\,
  (2\pi)^d2\omega_{\vec p_2}\delta^d(\vec q_1-\vec p_2)
\\
&\qquad+
  {}_{\mathrm{out}}\langle
  \vec q_1,\vec q_2
  |\vec p_1,\vec p_2
  \rangle_{\mathrm{in,conn}} .
\end{aligned}
\]

\section{Invariant Amplitudes}

The connected kernel contains the overall conservation law
\[
  q_1+\cdots+q_m=p_1+\cdots+p_n.
\]
It is conventional to define the invariant amplitude \(\mathcal M\) by
\[
\begin{aligned}
&{}_{\mathrm{out}}\langle
  \vec q_1,\ldots,\vec q_m
  |\vec p_1,\ldots,\vec p_n
  \rangle_{\mathrm{in,conn}}
\\
&\quad =
  \ii(2\pi)^D
  \delta^D(q_1+\cdots+q_m-p_1-\cdots-p_n)\,
  \mathcal M(q_1,\ldots,q_m|p_1,\ldots,p_n).
\end{aligned}
\]
With this convention the operator identity \(S=\mathbf 1+\ii T\) gives
\[
  {}_{\mathrm{out}}\langle q|T|p\rangle
  =
  (2\pi)^D\delta^D(P_{\mathrm{out}}-P_{\mathrm{in}})\,
  \mathcal M(q|p).
\]
Here \(P_{\mathrm{out}}=\sum_i q_i\) and
\(P_{\mathrm{in}}=\sum_j p_j\).

\begin{center}
\begin{tikzpicture}[scale=0.9,>=Stealth]
  \begin{scope}
    \draw[purple!80!black,very thick,->] (-2.7,1.0)--(-1.25,0.45);
    \draw[purple!80!black,very thick,->] (-2.7,-1.0)--(-1.25,-0.45);
    \draw[purple!80!black,very thick,->] (1.25,0.45)--(2.7,1.0);
    \draw[purple!80!black,very thick,->] (1.25,-0.45)--(2.7,-1.0);
    \node[circle,draw=black,thick,minimum size=1.45cm,
      fill=blue!7] at (0,0) {\(\widetilde G_c\)};
    \node[left] at (-2.75,1.0) {$p_1$};
    \node[left] at (-2.75,-1.0) {$p_2$};
    \node[right] at (2.75,1.0) {$q_1$};
    \node[right] at (2.75,-1.0) {$q_2$};
  \end{scope}
  \node at (3.45,0) {$\xrightarrow{\;\prod Z^{-1/2}\ii(k^2+m^2)\;}$};
  \begin{scope}[xshift=7.0cm]
    \draw[purple!80!black,very thick,->] (-1.65,0.8)--(-0.55,0.28);
    \draw[purple!80!black,very thick,->] (-1.65,-0.8)--(-0.55,-0.28);
    \draw[purple!80!black,very thick,->] (0.55,0.28)--(1.65,0.8);
    \draw[purple!80!black,very thick,->] (0.55,-0.28)--(1.65,-0.8);
    \node[circle,draw=black,thick,minimum size=1.15cm,
      fill=green!8] at (0,0) {\(\mathcal M\)};
  \end{scope}
\end{tikzpicture}
\end{center}

The diagram displays the operation. The left object is the connected
time-ordered Green function with its external one-particle poles. The right
object is the on-shell invariant amplitude obtained after the pole residues
have been extracted.

\section{Perturbative Evaluation After Reduction}

Perturbation theory supplies an expansion of the time-ordered Green function.
LSZ then turns that expansion into an expansion of matrix elements of the
scattering operator. This is the point where Feynman graphs acquire their
scattering-amplitude interpretation.

For the scalar interaction
\[
  \mathcal L_{\mathrm{int}}=-\frac{g}{4!}\phi^4,
\]
the leading connected four-point Green function is
\[
\begin{aligned}
&\widetilde G^{\mathrm{conn}}_4
  (q_1,q_2,-p_1,-p_2)
\\
&\quad =
  (2\pi)^D\delta^D(q_1+q_2-p_1-p_2)
  (-\ii g)
  \prod_{a=1}^4
  \frac{-\ii}{k_a^2+m^2-\ii0}
  +O(g^2),
\end{aligned}
\]
with
\[
  k_1=q_1,\quad k_2=q_2,\quad k_3=-p_1,\quad k_4=-p_2,
  \qquad
  Z=1+O(g^2).
\]
Applying LSZ gives
\[
  {}_{\mathrm{out}}\langle
  \vec q_1,\vec q_2
  |\vec p_1,\vec p_2
  \rangle_{\mathrm{in,conn}}
  =
  (2\pi)^D\delta^D(q_1+q_2-p_1-p_2)
  \left[-\ii g+O(g^2)\right].
\]
Therefore
\[
  \mathcal M(q_1,q_2|p_1,p_2)
  =
  -g+O(g^2)
\]
with the amplitude convention above. Higher orders are obtained from the
same Green-function expansion together with the external residue factors
\(Z^{-1/2}\).

\section{Scope of the Formula}

The LSZ theorem in this chapter applies to stable massive particles
represented by isolated one-particle poles in local-field Green functions. It
is a statement about boundary values of distributions after wave-packet
smearing. Its hypotheses include the existence of the asymptotic states
constructed in Haag--Ruelle theory and the existence of the relevant
time-ordered products.

Massless gauge theories, infraparticle sectors, confined degrees of freedom,
unstable resonances, and conformal theories require different asymptotic
objects or different observables. The common principle remains the same:
first construct the asymptotic Hilbert-space or detector data, then relate
that data to correlation functions in the regime where a reduction theorem is
available.
